% Source: https://artofproblemsolving.com/community/c3h1949089_2002_amc_12a_problem_17_potw_aops_wiki
% Stable identity: aops-topic-1949089-q01
\textbf{Lower bound.} The digits $4$, $6$ and $8$ cannot be the units digit of a prime: a number ending in an even digit is even, and $4,6,8$ themselves are not prime. So each of $4,6,8$ appears in the tens place or higher, contributing at least $40+60+80=180$ to the sum. The other six digits $1,2,3,5,7,9$ contribute at least their face values, $1+2+3+5+7+9=27$. Hence every such set has sum at least $180+27=207$.

\textbf{The bound is attained.} We need $4,6,8$ exactly in tens places and $1,2,3,5,7,9$ in units places (with no hundreds digits). For example
\[
\{2,\ 3,\ 5,\ 41,\ 67,\ 89\}
\]
works: $41$, $67$ and $89$ are prime, $2,3,5$ are prime, every nonzero digit is used once, and the sum is $2+3+5+41+67+89=207$.

So the smallest possible sum is $\boxed{\textbf{(B) }207}$. (Without searching blindly: once you know the tens digits must be $4,6,8$, you only have to pair them with units digits from $1,3,7,9$ to make primes, leaving the remaining digits as one-digit primes; $41,67,89$ is one quick choice.)
